\chapter{The linear cocycle and Furstenberg--Kesten}
\label{chap:cocycle}

\section{The measure-preserving system and the linear cocycle}

The Oseledets theorem studies the long-term growth of products of matrices driven by a
dynamical system. The data are a measure-preserving self-map \(T : X \to X\) of a
probability space \((X,\mu)\) and a matrix-valued generator
\(A : X \to \Matrix(\Fin d)(\Fin d)\,\R\). The fundamental object is the
\emph{iterated linear cocycle}: the product of the generator evaluated along the orbit
of \(x\). We use the convention that the newest factor sits on the left.

\begin{definition}[The entrywise measurable structure on matrices]\label{instMeasurableSpaceMatrix}
  \lean{Oseledets.instMeasurableSpaceMatrix}\leanok
  A matrix is a function \(m \to n \to \alpha\), and we equip
  \(\Matrix\, m\, n\, \alpha\) with the Pi (product) \(\sigma\)-algebra induced from the
  entry type \(\alpha\). For finitely many entries over a Borel \(\alpha\) this agrees
  with the Borel \(\sigma\)-algebra. Mathlib does not register this automatically because
  \(\Matrix\) is a \texttt{def} rather than reducibly the underlying Pi type, so the
  ambient Pi instance does not transfer; we install it explicitly.
\end{definition}

\begin{definition}[The iterated linear cocycle]\label{cocycle}
  \lean{Oseledets.cocycle}\leanok
  Given \(A : X \to \Matrix(\Fin d)(\Fin d)\,\R\) and \(T : X \to X\), define
  \(A^{(n)} = \texttt{cocycle}\,A\,T\,n\) by recursion on \(n\):
  \[
    A^{(0)}(x) = 1, \qquad A^{(n+1)}(x) = A^{(n)}(Tx)\cdot A(x).
  \]
  Unfolding, \(A^{(n)}(x) = A(T^{n-1}x)\cdots A(Tx)\,A(x)\), with the newest factor on
  the left. The matrix norm throughout is the scoped L2 operator norm, which is
  submultiplicative; vectors live in \(\mathrm{EuclideanSpace}\;\R\,(\Fin d)\) and the
  action is via \(\texttt{toEuclideanCLM}\).
\end{definition}

\begin{theorem}[The cocycle identity]\label{cocycle_add}
  \lean{Oseledets.cocycle_add}\leanok
  For all \(m, n \in \N\) and \(x \in X\),
  \[
    A^{(m+n)}(x) = A^{(m)}(T^{n}x)\cdot A^{(n)}(x).
  \]
\end{theorem}
\begin{proof}\leanok
  \uses{cocycle}
  Induction on \(n\) with \(x\) generalized. For \(n = 0\) both sides equal \(A^{(m)}(x)\).
  For the step, reassociate \(m + (n+1)\) and apply the defining recursion
  \(A^{(k+1)}(x) = A^{(k)}(Tx)\cdot A(x)\) on both sides, then use the inductive
  hypothesis at \(Tx\) together with \(T^{n+1}x = T^{n}(Tx)\) and associativity of matrix
  multiplication.
\end{proof}

\begin{definition}[One-sided log-integrability of the generator]\label{IntegrableLogNorm}
  \lean{Oseledets.IntegrableLogNorm}\leanok
  The hypothesis \(\texttt{IntegrableLogNorm}\,A\,\mu\) asserts that the positive part of
  the log-norm of the generator is integrable: \(\log^{+}\norm{A(\cdot)} \in L^{1}(\mu)\),
  where \(\log^{+} t = \max(\log t, 0)\). This is the standard integrability assumption of
  the Furstenberg--Kesten and Oseledets theorems; combined with the same hypothesis for
  the inverse generator \(A^{-1}\) it pins both extremal Lyapunov exponents in \(\R\).
\end{definition}

\begin{theorem}[Measurability of the cocycle iterates]\label{measurable_cocycle}
  \lean{Oseledets.measurable_cocycle}\leanok
  If \(A\) and \(T\) are measurable then for each \(n\) the iterate
  \(x \mapsto A^{(n)}(x)\) is measurable for the entrywise structure
  \ref{instMeasurableSpaceMatrix}.
\end{theorem}
\begin{proof}\leanok
  \uses{cocycle, instMeasurableSpaceMatrix}
  Induction on \(n\). The base case is a constant map. For the step, the recursion writes
  \(A^{(n+1)}\) as the product \((A^{(n)}\circ T)\cdot A\); matrix multiplication is jointly
  measurable on the Pi structure because each entry of a product is the finite sum
  \(\sum_k M_{ik}N_{kj}\) of products of measurable coordinate projections, so the result
  follows by composing the inductive hypothesis with \(T\) and multiplying by \(A\).
\end{proof}

\section{Measurability of the operator norm and the inverse}

To feed the cocycle into the ergodic theory we must know that the operator norm and the
matrix inverse are measurable on the entrywise structure. The subtlety is that Mathlib's
\(\texttt{Measurable.norm}\) is stated for a \texttt{BorelSpace}, whereas our matrix
\(\sigma\)-algebra is the Pi structure; the two coincide here because the L2
operator-norm topology is definitionally the Pi product topology.

\begin{lemma}[The Pi structure is an opens-measurable space]\label{instOpensMeasurableSpaceMatrix}
  \lean{Oseledets.instOpensMeasurableSpaceMatrix}\leanok
  The entrywise (Pi) measurable structure on \(\Matrix(\Fin d)(\Fin d)\,\R\) is an
  \texttt{OpensMeasurableSpace} for the L2 operator-norm topology, since that topology is
  installed (via \texttt{replaceTopology}) to be definitionally the Pi product topology, of
  which the Pi \(\sigma\)-algebra is exactly the Borel structure.
\end{lemma}

\begin{theorem}[Measurability of the L2 operator norm]\label{measurable_l2_opNorm}
  \lean{Oseledets.measurable_l2_opNorm}\leanok
  The map \(M \mapsto \norm{M}\) on \(\Matrix(\Fin d)(\Fin d)\,\R\) is measurable.
\end{theorem}
\begin{proof}\leanok
  \uses{instOpensMeasurableSpaceMatrix, instMeasurableSpaceMatrix}
  The norm is continuous for the operator-norm topology, and by
  \ref{instOpensMeasurableSpaceMatrix} that topology's Borel structure is the entrywise
  Pi structure; continuous maps into a Borel codomain are measurable.
\end{proof}

\begin{theorem}[Measurability of the determinant]\label{measurable_det}
  \lean{Oseledets.measurable_det}\leanok
  The determinant \(M \mapsto \det M\) is measurable.
\end{theorem}
\begin{proof}\leanok
  \uses{instMeasurableSpaceMatrix}
  By the Leibniz formula \(\det M = \sum_{\sigma}\operatorname{sgn}(\sigma)\prod_i M_{i,\sigma(i)}\),
  the determinant is a finite sum of finite products of measurable coordinate
  projections, hence measurable.
\end{proof}

\begin{theorem}[Measurability of the matrix inverse]\label{measurable_inv_matrix}
  \lean{Oseledets.measurable_inv_matrix}\leanok
  The inverse \(M \mapsto M^{-1}\) is measurable on the entrywise structure.
\end{theorem}
\begin{proof}\leanok
  \uses{measurable_det, instMeasurableSpaceMatrix}
  Writing \(M^{-1} = (\det M)^{-1}\cdot \operatorname{adj}(M)\), each entry is a ratio of
  polynomials in the entries. The adjugate is measurable (each entry is a determinant of a
  row update, again a polynomial in the entries), and \((\det M)^{-1}\) is measurable by
  \ref{measurable_det} and measurability of inversion on \(\R\); the product is therefore
  measurable entrywise.
\end{proof}

\section{Positivity and submultiplicativity of the log-norm}

To take logarithms cleanly we need the iterate norms to be strictly positive, which
forces the generator to be everywhere invertible (\(\det A \neq 0\)) and the dimension to
be nonzero. These hypotheses also make the \(\log\) of a product split as a genuine sum,
which is what yields subadditivity rather than a mere inequality with junk values.

\begin{lemma}[Invertibility of the iterates]\label{det_cocycle_ne_zero}
  \lean{Oseledets.det_cocycle_ne_zero}\leanok
  If \(\det A(x) \neq 0\) for every \(x\), then \(\det A^{(n)}(x) \neq 0\) for all
  \(n, x\).
\end{lemma}
\begin{proof}\leanok
  \uses{cocycle}
  Induction on \(n\). The base case is \(\det 1 = 1\). For the step,
  \(\det(A^{(n)}(Tx)\cdot A(x)) = \det A^{(n)}(Tx)\cdot \det A(x)\), and both factors are
  nonzero by the inductive hypothesis and the assumption on \(A\).
\end{proof}

\begin{lemma}[The unit matrix has norm one]\label{norm_one_matrix}
  \lean{Oseledets.norm_one_matrix}\leanok
  When \(d \neq 0\), \(\norm{(1 : \Matrix(\Fin d)(\Fin d)\,\R)} = 1\) for the L2 operator
  norm.
\end{lemma}
\begin{proof}\leanok
  For \(d \neq 0\) the space \(\mathrm{EuclideanSpace}\;\R\,(\Fin d)\) is nontrivial. The
  star-algebra equivalence \(\texttt{toEuclideanCLM}\) sends the identity matrix to the
  identity operator and preserves the norm, and the operator norm of the identity on a
  nontrivial space is \(1\). There is no \texttt{NormOneClass} instance for the matrix
  operator norm, so this must be proved by hand; note that at \(d = 0\) the statement is
  false (\(\norm{1} = 0\)), which is why \(d \neq 0\) is required.
\end{proof}

\begin{lemma}[Positivity of the iterate norms]\label{norm_cocycle_pos}
  \lean{Oseledets.norm_cocycle_pos}\leanok
  Assume \(\det A(x) \neq 0\) for all \(x\) and \(d \neq 0\). Then
  \(0 < \norm{A^{(n)}(x)}\) for every \(n, x\). The analogous statement
  \(\texttt{norm\_inv\_cocycle\_pos}\) holds for the inverse iterates
  \(\norm{(A^{(n)}(x))^{-1}}\).
\end{lemma}
\begin{proof}\leanok
  \uses{det_cocycle_ne_zero}
  If \(\norm{A^{(n)}(x)} = 0\) then \(A^{(n)}(x)\) is the zero matrix, whose determinant
  vanishes (as \(d \neq 0\)), contradicting \ref{det_cocycle_ne_zero}. The inverse case is
  identical using \(\det\bigl((A^{(n)}(x))^{-1}\bigr) = (\det A^{(n)}(x))^{-1} \neq 0\).
\end{proof}

\begin{theorem}[Subadditivity of the log-norm cocycle]\label{isSubadditiveCocycle_logNorm}
  \lean{Oseledets.isSubadditiveCocycle_logNorm}\leanok
  Assume \(\det A \neq 0\) everywhere and \(d \neq 0\). Then
  \(g_n(x) = \log\norm{A^{(n)}(x)}\) is a subadditive cocycle over \(T\):
  \[
    g_{m+n}(x) \le g_m(x) + g_n(T^{m}x).
  \]
\end{theorem}
\begin{proof}\leanok
  \uses{cocycle_add, norm_cocycle_pos}
  Rewrite \(m+n\) as \(n+m\) and apply the cocycle identity \ref{cocycle_add} to get
  \(A^{(m+n)}(x) = A^{(n)}(T^{m}x)\cdot A^{(m)}(x)\). Submultiplicativity of the operator
  norm and monotonicity of \(\log\) give
  \(\log\norm{A^{(m+n)}(x)} \le \log\bigl(\norm{A^{(n)}(T^{m}x)}\cdot\norm{A^{(m)}(x)}\bigr)\);
  both factors are strictly positive by \ref{norm_cocycle_pos}, so \(\log\) of the product
  splits as the sum \(\log\norm{A^{(n)}(T^{m}x)} + \log\norm{A^{(m)}(x)}\), which is the
  required bound after commuting the summands.
\end{proof}

\begin{theorem}[Subadditivity of the inverse log-norm cocycle]\label{isSubadditiveCocycle_logNorm_inv}
  \lean{Oseledets.isSubadditiveCocycle_logNorm_inv}\leanok
  Under the same hypotheses, \(g_n(x) = \log\norm{(A^{(n)}(x))^{-1}}\) is a subadditive
  cocycle over \(T\).
\end{theorem}
\begin{proof}\leanok
  \uses{cocycle_add, norm_cocycle_pos}
  As above, after \(\ref{cocycle_add}\) the inverse of the product reverses the order:
  \((A^{(m+n)}(x))^{-1} = (A^{(m)}(x))^{-1}\cdot (A^{(n)}(T^{m}x))^{-1}\). Submultiplicativity,
  monotone \(\log\), strict positivity of each inverse norm (\ref{norm_cocycle_pos}), and
  \(\log\) of a product splitting as a sum give the subadditive inequality.
\end{proof}

\section{Birkhoff-sum sandwich bounds}

The two-sided integrability hypothesis controls the log-norm by Birkhoff sums of
\(\log^{+}\norm{A}\) and \(\log^{+}\norm{A^{-1}}\) on both sides. These bounds drive both
the integrability of each level and the bounded-below proviso of Kingman's theorem.

\begin{lemma}[Upper Fekete bound]\label{logNorm_cocycle_le_birkhoffSum}
  \lean{Oseledets.logNorm_cocycle_le_birkhoffSum}\leanok
  For \(\det A \neq 0\) everywhere and \(d \neq 0\),
  \[
    \log\norm{A^{(n)}(x)} \le \sum_{k<n}\log^{+}\norm{A(T^{k}x)}.
  \]
\end{lemma}
\begin{proof}\leanok
  \uses{cocycle, norm_cocycle_pos}
  Induction on \(n\). The base case is \(0 \le 0\). For the step, the recursion gives
  \(A^{(n+1)}(x) = A^{(n)}(Tx)\cdot A(x)\); submultiplicativity and \(\log\)-splitting bound
  \(\log\norm{A^{(n+1)}(x)}\) by \(\log\norm{A(x)} + \log\norm{A^{(n)}(Tx)}\), and then
  \(\log\norm{A(x)} \le \log^{+}\norm{A(x)}\) together with the inductive hypothesis at
  \(Tx\) peels off one Birkhoff term.
\end{proof}

\begin{lemma}[Lower bound via the inverse Birkhoff sum]\label{neg_birkhoffSum_le_logNorm_cocycle}
  \lean{Oseledets.neg_birkhoffSum_le_logNorm_cocycle}\leanok
  Under the same hypotheses,
  \[
    -\sum_{k<n}\log^{+}\norm{(A(T^{k}x))^{-1}} \le \log\norm{A^{(n)}(x)}.
  \]
\end{lemma}
\begin{proof}\leanok
  \uses{logNorm_inv_cocycle_le_birkhoffSum, norm_one_matrix, det_cocycle_ne_zero}
  From \(A^{(n)}(x)\cdot (A^{(n)}(x))^{-1} = 1\) and submultiplicativity,
  \(1 = \norm{1} \le \norm{A^{(n)}(x)}\cdot\norm{(A^{(n)}(x))^{-1}}\) (using
  \ref{norm_one_matrix} and \ref{det_cocycle_ne_zero}), so taking logs gives
  \(0 \le \log\norm{A^{(n)}(x)} + \log\norm{(A^{(n)}(x))^{-1}}\). Bounding the inverse
  log-norm above by its Birkhoff sum \ref{logNorm_inv_cocycle_le_birkhoffSum} and
  rearranging yields the claim.
\end{proof}

\begin{lemma}[Upper Fekete bound for the inverse cocycle]\label{logNorm_inv_cocycle_le_birkhoffSum}
  \lean{Oseledets.logNorm_inv_cocycle_le_birkhoffSum}\leanok
  Under the same hypotheses,
  \[
    \log\norm{(A^{(n)}(x))^{-1}} \le \sum_{k<n}\log^{+}\norm{(A(T^{k}x))^{-1}}.
  \]
\end{lemma}
\begin{proof}\leanok
  \uses{cocycle, norm_cocycle_pos}
  Induction on \(n\), exactly as for the forward bound but with the inverse generator: the
  recursion and \(\texttt{mul\_inv\_rev}\) give
  \((A^{(n+1)}(x))^{-1} = (A(x))^{-1}\cdot (A^{(n)}(Tx))^{-1}\), and submultiplicativity plus
  \(\log\)-splitting and \(\log\norm{(A(x))^{-1}} \le \log^{+}\norm{(A(x))^{-1}}\) peel off
  one Birkhoff term.
\end{proof}

\section{Integrability of each level}

\begin{lemma}[Integral of a Birkhoff sum]\label{integral_birkhoffSum}
  \lean{Oseledets.integral_birkhoffSum}\leanok
  For measure-preserving \(T\) and integrable \(f\),
  \(\int \sum_{k<n} f(T^{k}x)\,d\mu = n\int f\,d\mu\).
\end{lemma}
\begin{proof}\leanok
  Each composition \(f\circ T^{k}\) is integrable and integral-preserving because
  \(T^{k}\) is measure-preserving, so \(\int f\circ T^{k}\,d\mu = \int f\,d\mu\). Summing
  the \(n\) equal terms over the finite range gives \(n\int f\,d\mu\).
\end{proof}

\begin{theorem}[Integrability of the log-norm levels]\label{integrable_logNorm_cocycle}
  \lean{Oseledets.integrable_logNorm_cocycle}\leanok
  Let \(T\) be measure-preserving for a finite measure \(\mu\), \(A\) measurable and
  everywhere invertible, \(d \neq 0\), with both \(\log^{+}\norm{A}\) and
  \(\log^{+}\norm{A^{-1}}\) integrable. Then each \(x \mapsto \log\norm{A^{(n)}(x)}\) is
  integrable. The companion \(\texttt{integrable\_logNorm\_inv\_cocycle}\) gives the same
  for the inverse iterates.
\end{theorem}
\begin{proof}\leanok
  \uses{logNorm_cocycle_le_birkhoffSum, neg_birkhoffSum_le_logNorm_cocycle, integral_birkhoffSum, measurable_l2_opNorm, IntegrableLogNorm}
  The level \(g_n\) is sandwiched between \(-B_n^{-}\) and \(B_n^{+}\), where
  \(B_n^{\pm}\) are the Birkhoff sums of \(\log^{+}\norm{A}\) and \(\log^{+}\norm{A^{-1}}\)
  (Lemmas \ref{logNorm_cocycle_le_birkhoffSum}, \ref{neg_birkhoffSum_le_logNorm_cocycle}).
  Both \(B_n^{\pm}\) are nonnegative and integrable (finite sums of integrable, m.p.
  precompositions), so \(\abs{g_n} \le B_n^{+} + B_n^{-}\) pointwise. Since \(g_n\) is
  measurable (\ref{measurable_l2_opNorm} composed with the cocycle), domination by the
  integrable \(B_n^{+}+B_n^{-}\) gives integrability.
\end{proof}

\section{The Furstenberg--Kesten theorems}

Feeding the subadditive cocycles into Kingman's ergodic theorem produces the two extremal
Lyapunov exponents as a.e.-constant limits. The bounded-below proviso of Kingman is
supplied by the cross integrability hypothesis: for the top exponent the lower bound comes
from \(\log^{+}\norm{A^{-1}}\), and vice versa. This is precisely where the second
integrability hypothesis keeps the limit finite in \(\R\) rather than \(-\infty\).

\begin{theorem}[Furstenberg--Kesten, top exponent]\label{furstenbergKesten_norm}
  \lean{Oseledets.furstenbergKesten_norm}\leanok
  Let \(T\) be ergodic for a probability measure \(\mu\), let \(A\) be measurable and
  everywhere invertible with \(\log^{+}\norm{A}, \log^{+}\norm{A^{-1}} \in L^{1}(\mu)\).
  Then there is a constant \(\lambda_1 \in \R\) (the top Lyapunov exponent) with
  \[
    \frac{1}{n}\log\norm{A^{(n)}(x)} \xrightarrow{n\to\infty} \lambda_1
    \qquad \text{for } \mu\text{-a.e. } x.
  \]
\end{theorem}
\begin{proof}\leanok
  \uses{isSubadditiveCocycle_logNorm, integrable_logNorm_cocycle, neg_birkhoffSum_le_logNorm_cocycle, integral_birkhoffSum, tendsto_kingman_ergodic, lambdaBar}
  If \(d = 0\) the matrix algebra is trivial, every norm is \(0\), and the limit is the
  constant \(0\). Otherwise set \(d \neq 0\) and let \(g_n = \log\norm{A^{(n)}}\). It is a
  subadditive cocycle (\ref{isSubadditiveCocycle_logNorm}) with each level integrable
  (\ref{integrable_logNorm_cocycle}). For the bounded-below proviso, the lower bound
  \ref{neg_birkhoffSum_le_logNorm_cocycle} and the Birkhoff integral identity
  \ref{integral_birkhoffSum} give \((\int g_{n+1})/(n+1) \ge -\int \log^{+}\norm{A^{-1}}\,d\mu\),
  a constant lower bound. The ergodic Kingman theorem \ref{tendsto_kingman_ergodic} then
  returns the a.e. constant limit, which is the top exponent \(\lambda_1\) (see
  \ref{lambdaBar}).
\end{proof}

\begin{theorem}[Furstenberg--Kesten, bottom exponent]\label{furstenbergKesten_norm_inv}
  \lean{Oseledets.furstenbergKesten_norm_inv}\leanok
  Under the same hypotheses there is a constant \(\lambda \in \R\) with
  \[
    \frac{1}{n}\log\norm{(A^{(n)}(x))^{-1}} \xrightarrow{n\to\infty} \lambda
    \qquad \text{for } \mu\text{-a.e. } x,
  \]
  so the bottom Lyapunov exponent \(\lambda_k = -\lambda\) exists and is finite.
\end{theorem}
\begin{proof}\leanok
  \uses{isSubadditiveCocycle_logNorm_inv, integrable_logNorm_cocycle, logNorm_cocycle_le_birkhoffSum, integral_birkhoffSum, norm_one_matrix, tendsto_kingman_ergodic}
  Identical to the top case with the inverse subadditive cocycle
  \(g_n = \log\norm{(A^{(n)})^{-1}}\) (\ref{isSubadditiveCocycle_logNorm_inv}). The roles of
  the two integrability hypotheses swap: the bounded-below proviso now uses the lower
  bound \(-\sum_{k<n}\log^{+}\norm{A(T^{k}x)} \le \log\norm{(A^{(n)}(x))^{-1}}\), obtained
  from \(\norm{1} \le \norm{A^{(n)}}\cdot\norm{(A^{(n)})^{-1}}\) (\ref{norm_one_matrix}) and
  the forward Fekete bound \ref{logNorm_cocycle_le_birkhoffSum}, together with the integral
  identity \ref{integral_birkhoffSum}. Ergodic Kingman \ref{tendsto_kingman_ergodic}
  delivers the a.e. constant limit.
\end{proof}
